\documentclass[10pt,a4paper]{amsart}
\usepackage{amsmath,amssymb,amsthm,graphicx}
% For commutative diagrams we use Xy-pic
\usepackage[all]{xy}

% We use 2cell for 2-commutative diagrams.
\xyoption{2cell}
% Theorem environments.
%
\theoremstyle{plain}
\newtheorem{theorem}[subsection]{Theorem}
\newtheorem{proposition}[subsection]{Proposition}
\newtheorem{lemma}[subsection]{Lemma}

\theoremstyle{definition}
\newtheorem{definition}[subsection]{Definition}
\newtheorem{example}[subsection]{Example}
\newtheorem{exercise}[subsection]{Exercise}
\newtheorem{situation}[subsection]{Situation}

\theoremstyle{remark}
\newtheorem{remark}[subsection]{Remark}
\newtheorem{remarks}[subsection]{Remarks}

\numberwithin{equation}{subsection}

% Macros
%
\def\lim{\mathop{\mathrm{lim}}\nolimits}
\def\colim{\mathop{\mathrm{colim}}\nolimits}
\def\Spec{\mathop{\mathrm{Spec}}}
\def\Hom{\mathop{\mathrm{Hom}}\nolimits}
\def\Ext{\mathop{\mathrm{Ext}}\nolimits}
\def\SheafHom{\mathop{\mathcal{H}\!\mathit{om}}\nolimits}
\def\SheafExt{\mathop{\mathcal{E}\!\mathit{xt}}\nolimits}
\def\Sch{\mathit{Sch}}
\def\Mor{\mathop{\mathrm{Mor}}\nolimits}
\def\Ob{\mathop{\mathrm{Ob}}\nolimits}
\def\Sh{\mathop{\mathit{Sh}}\nolimits}
\def\NL{\mathop{N\!L}\nolimits}
\def\CH{\mathop{\mathrm{CH}}\nolimits}
\def\proetale{{pro\text{-}\acute{e}tale}}
\def\etale{{\acute{e}tale}}
\def\QCoh{\mathit{QCoh}}
\def\Ker{\mathop{\mathrm{Ker}}}
\def\Im{\mathop{\mathrm{Im}}}
\def\Coker{\mathop{\mathrm{Coker}}}
\def\Coim{\mathop{\mathrm{Coim}}}

% Boxtimes
%
\DeclareMathSymbol{\boxtimes}{\mathbin}{AMSa}{"02}

%
% Macros for moduli stacks/spaces
%
\def\QCohstack{\mathcal{QC}\!\mathit{oh}}
\def\Cohstack{\mathcal{C}\!\mathit{oh}}
\def\Spacesstack{\mathcal{S}\!\mathit{paces}}
\def\Quotfunctor{\mathrm{Quot}}
\def\Hilbfunctor{\mathrm{Hilb}}
\def\Curvesstack{\mathcal{C}\!\mathit{urves}}
\def\Polarizedstack{\mathcal{P}\!\mathit{olarized}}
\def\Complexesstack{\mathcal{C}\!\mathit{omplexes}}
% \Pic is the operator that assigns to X its picard group, usage \Pic(X)
% \Picardstack_{X/B} denotes the Picard stack of X over B
% \Picardfunctor_{X/B} denotes the Picard functor of X over B
\def\Pic{\mathop{\mathrm{Pic}}\nolimits}
\def\Picardstack{\mathcal{P}\!\mathit{ic}}
\def\Picardfunctor{\mathrm{Pic}}
\def\Deformationcategory{\mathcal{D}\!\mathit{ef}}
\usepackage{cleveref}
\crefname{originalreference}{original reference}{original references}
\Crefname{originalreference}{Original reference}{Original references}
\makeatletter
\newlabel{lemma-what-integral}{{04C2}{source}{}{sourcepointer-1}{}}
\newlabel{lemma-what-integral@cref}{{[originalreference][0][]04C2}{[1][0][]source}{}{}{}}
\makeatother
\crefname{originalreference}{original reference}{original references}
\Crefname{originalreference}{Original reference}{Original references}
\makeatletter
\newlabel{remark-base-change-holds}{{0F02}{source}{}{sourcepointer-2}{}}
\newlabel{remark-base-change-holds@cref}{{[originalreference][0][]0F02}{[1][0][]source}{}{}{}}
\makeatother
\title{Base change in degree q descends along an integral base change of the target}
\author{The Stacks Project Authors}
\date{stacks-project:etale-cohomology.tex:local-snapshot-20261011; GNU FDL 1.2 or later}
\begin{document}
\maketitle
\noindent\textit{Editorial note.} All delivered mathematical text is the author's own, copied from the pinned Stacks Project chapter file etale-cohomology.tex; no statement or proof step was rewritten, reordered or removed.\par\medskip
\noindent\textit{Editorial note.} The counted proof is the author's complete argument: the proof is written out in full: the exactness of the pushforwards along integral morphisms, the two displayed rewritings through the base change identity, and the conclusion from the hypothesis.\par\medskip
\noindent\textit{Editorial note.} The label lemma-base-change-q-integral-bottom names the counted statement, and the proof environment immediately following it in the source is delivered verbatim as its proof.\par\medskip
\noindent\textit{Editorial note.} The counted statement is carried by the author's own lemma environment in the chapter source.\par\medskip
\noindent\textit{Editorial note.} The statement's two displayed diagrams are the author's and are delivered verbatim; the two cited inputs are kept as external pointers with their canonical Stacks tags.\par\medskip
\noindent\textit{Allowed standard prerequisites (editorial).}\begin{itemize}
\item Exactness of the pushforward along an integral morphism. Same chapter, the lemma stating that for an integral morphism the derived pushforward is the underived one. Supplies Rpi\_* = pi\_* and Rpi'\_* = pi'\_* and the fact that the pushforward reflects isomorphisms in the relevant sense.
\item Base change of the pushforward along a cartesian square. Same chapter, the base change identity (f')\textasciicircum{}\{-1\}pi\_* = pi'\_*(f'')\textasciicircum{}\{-1\} for the cartesian square in the statement. Supplies the two displayed rewritings of the two sides after applying pi'\_*.
\end{itemize}\medskip
\ifdefined\hypertarget\hypertarget{sourcepointer-1}{}\fi
\noindent\textit{Source reference.} Original reference 04C2: sources/stacks-etale-cohomology/pinned-original-source.tex, source lines 6056--6056. Prerequisite pointer invoked at proof line(s) 17035; it names a standard fact of the same chapter that the counted argument uses, and every step of the counted proof built on it is retained verbatim in the delivered body.\par\medskip
\ifdefined\hypertarget\hypertarget{sourcepointer-2}{}\fi
\noindent\textit{Source reference.} Original reference 0F02: sources/stacks-etale-cohomology/pinned-original-source.tex, source lines 16848--16848. Prerequisite pointer invoked at proof line(s) 16998; it names a standard fact of the same chapter that the counted argument uses, and every step of the counted proof built on it is retained verbatim in the delivered body.\par\medskip
\begin{lemma}
\label{lemma-base-change-q-integral-bottom}
With $f : X \to S$ and $n$ as in Remark \ref{remark-base-change-holds}
assume for some $q \geq 1$ we have $BC(f, n, q - 1)$. Consider
commutative diagrams
$$
\vcenter{
\xymatrix{
X \ar[d]_f &
X' \ar[d]_{f'} \ar[l] &
X'' \ar[l]^{\pi'} \ar[d]_{f''} &
Y \ar[l]^{h'} \ar[d]^e \\
S &
S' \ar[l] &
S'' \ar[l]_\pi &
T \ar[l]_{g'}
}
}
\quad\text{and}\quad
\vcenter{
\xymatrix{
X' \ar[d]_{f'} & & Y \ar[ll]^{h = h' \circ \pi'} \ar[d]^e \\
S' & & T \ar[ll]_{g = g' \circ \pi}
}
}
$$
where all squares are cartesian, $g'$ quasi-compact and quasi-separated, and
$\pi$ is integral. Let $\mathcal{F}$ be an abelian sheaf
on $T_\etale$ annihilated by $n$. If the base change map
$$
(f')^{-1}R^qg_*\mathcal{F} \longrightarrow R^qh_*e^{-1}\mathcal{F}
$$
is an isomorphism, then the base change map
$(f'')^{-1}R^qg'_*\mathcal{F} \to R^qh'_*e^{-1}\mathcal{F}$
is an isomorphism.
\end{lemma}

\begin{proof}
Since $\pi$ and $\pi'$ are integral we have $R\pi_* = \pi_*$ and
$R\pi'_* = \pi'_*$, see Lemma \ref{lemma-what-integral}.
We also have $(f')^{-1}\pi_* = \pi'_*(f'')^{-1}$. Thus we see that
$\pi'_*(f'')^{-1}R^qg'_*\mathcal{F} = (f')^{-1}R^qg_*\mathcal{F}$
and
$\pi'_*R^qh'_*e^{-1}\mathcal{F} = R^qh_*e^{-1}\mathcal{F}$.
Thus the assumption means that our map becomes an
isomorphism after applying the functor $\pi'_*$.
Hence we see that it is an isomorphism by Lemma \ref{lemma-what-integral}.
\end{proof}

\end{document}
